\chapter{Konfigurasi Lanjutan}
\label{ch:konfigurasi-lanjutan}

\section{Berkas Konfigurasi}

PakAI menyimpan konfigurasi di:
\begin{itemize}
  \item \texttt{\$XDG\_CONFIG\_HOME/pakai/config.toml}: jika variabel tersebut diset.
  \item \texttt{\textasciitilde/.config/pakai/config.toml}: lokasi default.
\end{itemize}

Berkas dibuat secara otomatis saat pertama kali menyimpan konfigurasi.

\section{Melihat Semua Konfigurasi}

\begin{lstlisting}[style=terminal]
$ pakai config list
\end{lstlisting}

\begin{figure}[h]
  \centering
  \includegraphics[width=0.6\linewidth]{pakai-config-list}
  \caption{Output \texttt{pakai config list}: semua konfigurasi aktif}
\end{figure}

\section{Mengubah Konfigurasi}

Gunakan \cmd{pakai config set \textit{kunci} \textit{nilai}} untuk mengubah nilai:

\begin{lstlisting}[style=terminal]
$ pakai config set daemon.port 7732
$ pakai config get daemon.port
\end{lstlisting}

\section{Referensi Kunci Konfigurasi}

\subsection{Konfigurasi Daemon}

\begin{center}
\begin{tabular}{lllp{5cm}}
\toprule
\textbf{Kunci} & \textbf{Tipe} & \textbf{Default} & \textbf{Keterangan} \\
\midrule
\texttt{daemon.port} & integer & 7731 & Port HTTP daemon (1--65535) \\
\texttt{daemon.poll\_interval} & integer & 30 & Interval polling dalam detik ($\geq 1$) \\
\bottomrule
\end{tabular}
\end{center}

\begin{lstlisting}[style=terminal]
$ pakai config set daemon.poll_interval 15
$ pakai config set daemon.port 7732
\end{lstlisting}

\subsection{Konfigurasi Tampilan}

\begin{center}
\begin{tabular}{lllp{5cm}}
\toprule
\textbf{Kunci} & \textbf{Tipe} & \textbf{Default} & \textbf{Keterangan} \\
\midrule
\texttt{display.separator} & string & \texttt{ | } & Pemisah antar penyedia \\
\texttt{display.order} & string & (auto) & Urutan penyedia, dipisah koma \\
\bottomrule
\end{tabular}
\end{center}

\begin{lstlisting}[style=terminal]
$ pakai config set display.separator " - "
$ pakai config set display.order "claude,opencode,openai"
\end{lstlisting}

\subsection{Konfigurasi Ambang Batas}

\begin{center}
\begin{tabular}{lllp{5cm}}
\toprule
\textbf{Kunci} & \textbf{Tipe} & \textbf{Default} & \textbf{Keterangan} \\
\midrule
\texttt{thresholds.warning} & integer & 50 & Ambang peringatan (0--100\%) \\
\texttt{thresholds.critical} & integer & 80 & Ambang kritis (0--100\%) \\
\bottomrule
\end{tabular}
\end{center}

\begin{lstlisting}[style=terminal]
$ pakai config set thresholds.warning 60
$ pakai config set thresholds.critical 85
\end{lstlisting}

\subsection{Konfigurasi Per Penyedia}

Setiap penyedia dapat dikonfigurasi dengan format \texttt{provider.\textit{id}.\textit{kunci}}:

\begin{center}
\begin{tabular}{lllp{4.5cm}}
\toprule
\textbf{Kunci} & \textbf{Tipe} & \textbf{Default} & \textbf{Keterangan} \\
\midrule
\texttt{provider.\textit{id}.enabled} & bool & true & Aktifkan/nonaktifkan penyedia \\
\texttt{provider.\textit{id}.label} & string & (id) & Label tampilan kustom \\
\texttt{provider.\textit{id}.limit} & float & 0 & Batas biaya dalam USD ($\geq 0$) \\
\bottomrule
\end{tabular}
\end{center}

ID penyedia yang tersedia: \texttt{claude}, \texttt{opencode}, \texttt{openai},
\texttt{opencode/openai}, \texttt{opencode/opencode-go}.

\begin{lstlisting}[style=terminal]
# Atur label kustom untuk Claude
$ pakai config set provider.claude.label "Klaude"

# Atur batas bulanan untuk OpenCode Go
$ pakai config set provider.opencode-go.limit 60

# Nonaktifkan penyedia OpenAI sementara
$ pakai config set provider.openai.enabled false
\end{lstlisting}

\section{Contoh Berkas Konfigurasi Lengkap}

\begin{lstlisting}[style=toml]
[daemon]
port = 7731
poll_interval = 30

[display]
separator = " | "
order = "claude,opencode/opencode-go,openai"

[thresholds]
warning = 50
critical = 80

[provider.claude]
enabled = true
label = "Claude"

[provider.openai]
enabled = true
label = "Codex"

[provider.opencode]
enabled = true
label = "OpenCode"

[provider."opencode/opencode-go"]
enabled = true
limit = 60.0
\end{lstlisting}

\section{Mock Provider (untuk Pengujian)}

PakAI mendukung penyedia palsu (\textit{mock}) untuk keperluan pengujian tampilan:

\begin{lstlisting}[style=terminal]
# Buat mock dengan 75% penggunaan
$ pakai provider mock test-ai --used 75 --limit 100 --unit messages

# Hapus mock
$ pakai provider unmock test-ai
\end{lstlisting}
